\clearpage
\thispagestyle{empty}
\begin{rpgdark}[VERBETE DO ARQUIVO · A CRIATURA]
\centering
{\textgoth{\fontsize{26}{28}\selectfont A CRIATURA}}\\[0.25em]
{\scshape\bfseries\small ARQUIVO 1.260 · CLASSIFICAÇÃO: [ILEGÍVEL]}\\[0.3em]
{\color{rpggold}\rule{0.55\textwidth}{0.6pt}}
\end{rpgdark}

\vspace{0.5em}
\begin{rpgdark}[FICHA]
\small
\renewcommand{\arraystretch}{1.3}
\begin{tabularx}{\textwidth}{@{}lX@{}}
\textbf{Nome:} & [ilegível] \\
\textbf{Classe:} & consequência \\
\textbf{Nível:} & 1.260 \\
\textbf{Ataque:} & leitura \\
\textbf{Defesa:} & ser passada adiante \\
\textbf{Fraqueza:} & [vazio] \\
\textbf{Teste:} & 1d20 + Lucidez (CD 15) \\
\end{tabularx}
\end{rpgdark}

\vspace{0.5em}
\begin{rpgbox}[NOTA DO ARQUIVISTA]
\centering
\itshape ``Ela não é um monstro. Ela é uma consequência.''
\end{rpgbox}

\vspace{0.5em}
\begin{rpgbox}[REAÇÃO · role 1d6]
\small
\begin{tabularx}{\textwidth}{@{}cX@{}}
\toprule
\textbf{d6} & \textbf{O que acontece} \\
\midrule
1--2 & Silêncio. O cheiro passa. Não confie. \\
3--4 & Cheiro doce de terra molhada; uma página vira sozinha. \\
5 & O olho direito brilha. Teste de Lucidez para não seguir. \\
6 & A página seguinte está escrita com a sua letra. \\
\bottomrule
\end{tabularx}
\end{rpgbox}

\vspace{0.5em}
\noindent{\footnotesize\itshape Verbete ficcional; não é um stat block jogável. O arquivo não fecha. A criatura não tem ficha completa porque ninguém que a preencheu voltou para entregar.}
\clearpage
